\section{Glossary, notation and references}

\begin{multicols}{2}
\subsection{Compact glossary}
\footnotesize
\textbf{Additive RGB composition.} Combination of red, green and blue light; full coincident channels produce white on an emissive display.

\textbf{Antialiasing.} Partial edge coverage or supersampling used to reduce stair-stepping and flicker.

\textbf{Carrier.} Repeated line, lattice or path whose local width or size carries image information.

\textbf{Chromatic fringe.} Coloured boundary produced by relative RGB displacement.

\textbf{Coverage.} Fraction of a local cell occupied by one channel's geometry.

\textbf{Cycles per degree.} Number of repeated cycles subtending one degree of visual angle.

\textbf{Density.} Stripe count, dot spacing or spiral turn spacing, depending on the renderer.

\textbf{Duty cycle.} Fraction of a period occupied by a stripe.

\textbf{Halftone.} Continuous-tone representation using variable-size repeated dots.

\textbf{Point-spread function.} Spatial response to an ideal point, used here as an approximation of optical and neural blur.

\textbf{Spatial frequency.} Rate of repetition or change across visual space.

\textbf{Strength.} Amplitude of intensity-to-geometry modulation.

\subsection{Further reading}
\begin{thebibliography}{9}
\footnotesize
\bibitem{campbell1968} F. W. Campbell and J. G. Robson, ``Application of Fourier analysis to the visibility of gratings,'' \emph{J. Physiol.}, 197(3), 551--566, 1968. \href{https://doi.org/10.1113/jphysiol.1968.sp008574}{doi:10.1113/jphysiol.1968.sp008574}.

\bibitem{mullen1985} K. T. Mullen, ``The contrast sensitivity of human colour vision to red--green and blue--yellow chromatic gratings,'' \emph{J. Physiol.}, 359, 381--400, 1985. \href{https://doi.org/10.1113/jphysiol.1985.sp015591}{doi:10.1113/jphysiol.1985.sp015591}.

\bibitem{peli1990} E. Peli, ``Contrast in complex images,'' \emph{JOSA A}, 7(10), 2032--2040, 1990. \href{https://doi.org/10.1364/JOSAA.7.002032}{doi:10.1364/JOSAA.7.002032}.

\bibitem{devalois1988} R. L. De Valois and K. K. De Valois, \emph{Spatial Vision}. Oxford University Press, 1988. ISBN 0-19-505019-3.
\end{thebibliography}
\end{multicols}

\subsection{Principal notation}
\begin{table}[H]
\centering\footnotesize
\begin{tabularx}{\textwidth}{c >{\raggedright\arraybackslash}X c >{\raggedright\arraybackslash}X c >{\raggedright\arraybackslash}X}
\toprule
Symbol & Meaning & Symbol & Meaning & Symbol & Meaning \\
\midrule
$W,H$ & Output dimensions & $x,y$ & Output coordinates & $u,v$ & Source coordinates \\
$R,G,B$ & Source channels & $M_q$ & Channel mask & $a_q$ & Geometric amplitude \\
$P$ & Carrier spacing & $N$ & Stripe count & $d$ & Colour separation \\
$r$ & Dot/polar radius & $T$ & Turn spacing or duration & $F$ & Frames per second \\
$D$ & Viewing distance & $f_{\mathrm{cpd}}$ & Cycles per degree & $L_0$ & 1600 px reference size \\
\bottomrule
\end{tabularx}
\end{table}

\subsection{Document status}
This manual describes Percepta version 0.31.0 and the interface supplied in July 2026. The equations form a coherent technical specification of observed behaviour; they are not a verbatim listing of source-code operations. Default values and figures should be updated when the application changes.

\begin{tcolorbox}[colback=PerceptaNavy,colframe=PerceptaNavy,arc=2mm,boxrule=0pt,left=4mm,right=4mm,top=2.5mm,bottom=2.5mm]
\centering\color{white}
{\large\bfseries PERCEPTA}\quad Perceptual pattern image generator\par
\href{https://www.virgile-adam.com}{\color{PerceptaGreen}www.virgile-adam.com}
\quad|\quad
\href{mailto:virgile.adam@ibs.fr}{\color{PerceptaBlue}virgile.adam@ibs.fr}
\end{tcolorbox}
